The rational fast path terminates uniquely by recursion soundness.  If it returns identification, relabelling the terminal \(\operatorname{ID\_RAT}\) preserves its derivation, and recursion soundness proves \(F\in\mathscr R(\mathcal P)\) and equality with the causal query throughout the broad compatibility fiber.  A proved identically zero evidence terminal is likewise sound.  Every other local terminal invokes the canonical phase.

By exact global canonical frontier, nonempty \(\mathcal C(\mathcal P)\) implies that
\[
S=\Theta_G^{\mathrm{can}}(\mathcal P)
\]
is a nonempty compact semialgebraic set whose target-cell range equals the broad SCM range.  Its polynomial coefficients are rational except for the supplied values, which are effectively real algebraic.

We now give the terminating construction used by the fallback.  For polynomials \(p,q\), with \(q>0\) on \(S\), form
\[
\varphi_{p/q}(t)\Longleftrightarrow
\exists x\,[x\in S\ \wedge\ q(x)>0\ \wedge\ q(x)t-p(x)=0].
\tag{1}
\]
Apply effective real-algebraic quantifier elimination with the variable and polynomial order recorded in the certificate.  The established CAD result projects the finite polynomial list, isolates the real algebraic roots of the univariate projection polynomials, lifts sign-invariant cells, and supplies algebraic sample points.  It therefore returns an exact quantifier-free description of the compact image \(I_{p/q}\).  The least and greatest cells of this image are closed point cells; lifting them through (1) returns \(x_-,x_+\in S\) attaining the endpoints.  Recomputing projection polynomials, checking isolating intervals or Thom signs, checking every lifted sign table and constraint, and substituting the samples into (1) is a finite exact verifier.  Different legal CAD choices may produce different traces, but the established termination theorem applies to each and every accepted trace proves the same endpoint assertions.

For an unconditional target use \(p=\Phi_{T,\mathrm{can}}^\star\) and \(q=1\).  Let the attained extrema be \(L\le U\).  If \(L=U\), exact range equality makes the target constant and validates \(\operatorname{ID\_CAN}(L,\zeta)\).  Conversely, identification forces the extrema to agree.  If \(L<U\), lifted endpoint parameters canonically realize two SCM families on the original alphabets with the same menu and unequal target values, validating \(\operatorname{PAIR\_CAN}\).

For a conditional query, first apply the same construction to the nonnegative evidence polynomial \(D=\Phi_{E,\mathrm{can}}^\star\), obtaining \(0\le e_-^{\mathrm{can}}\le e_+^{\mathrm{can}}\).  If \(e_+^{\mathrm{can}}=0\), evidence vanishes throughout the fiber.  If \(e_-^{\mathrm{can}}=0<e_+^{\mathrm{can}}\), lifted endpoint samples give respectively a zero- and a positive-evidence original-alphabet model, so the causal conditional is not defined on the whole fiber.  These prove the first two conditional terminals without division.

In the remaining case \(e_-^{\mathrm{can}}>0\), set \(N=\Phi_{B_\rho\cap E,\mathrm{can}}^\star\) and apply (1) with \(p=N,q=D\).  The construction returns attained sharp ratio extrema \(L_Q\le U_Q\).  Exact canonical range equality identifies this ratio with \(P_{M^\star}(B_\rho\mid E)\).  Equality validates \(\operatorname{ID\_CAN}\); strict inequality yields two endpoint canonical realizations and validates \(\operatorname{PAIR\_CAN}\).  The evidence cases are mutually exclusive and exhaustive, as are equality and strict inequality of the ratio extrema.  Hence every legal finite run returns one valid certified output, with no uniqueness claim about its CAD representation.